\clearpage
\phantomsection
\refstepcounter{section}
\addcontentsline{toc}{section}{APÊNDICE \Alph{section} -- EXEMPLOS DE INTERAÇÃO MANUAL E DISCUSSÃO QUALITATIVA}
\begin{center}
\textbf{APÊNDICE \Alph{section} -- EXEMPLOS DE INTERAÇÃO MANUAL E DISCUSSÃO QUALITATIVA}
\end{center}
\label{sec:apendice-interacao}

Este apêndice reúne exemplos completos de interação manual utilizados para complementar a demonstração qualitativa apresentada na seção de experimentos. O objetivo é tornar auditáveis, em formato legível, um caso de bom desempenho em consulta temática delimitada, um caso em que o próprio sistema reconhece limites do contexto e evita atribuições indevidas de autoria, e um caso mais desafiador, no qual a pergunta exige consolidação legislativa ampla, articulação entre autores, cronograma, orçamento e consenso.

\subsection{Caso 1: resposta auditável em consulta temática delimitada}\label{sec:anexo-caso-bom}

\textbf{Pergunta}

Quais são os principais argumentos de Paulo Paim contra a capitalização da Previdência e que exemplos ele usa para sustentar essa posição?

\textbf{Resposta}

\textbf{Principais argumentos de Paulo Paim contra a capitalização da Previdência}

\begin{enumerate}
  \item \textbf{``Capitalização'' equivale, na prática, a ``poupança individual'', o que fragiliza a proteção previdenciária.}

  Ele afirma que ``regime de capitalização me preocupa muito'' porque ``significa uma poupança individual'', em que ``quem vai cuidar do meu dinheiro é aquele fundo'' e ``se o fundo aplicar mal, o resultado é zero'' em rendimento.

  \item \textbf{Perde-se a lógica solidária, e sobra apenas uma parcela do trabalhador.}

  Ele argumenta que se retiram diversas fontes de financiamento e se deixa ``só 10\% do empregado'', dizendo que ``claro que não vai dar'', pois ``tu tens 100, tira 90''.

  \item \textbf{Com esse modelo, a aposentadoria tende a ser muito baixa e pode não garantir continuidade para familiares.}

  Ele descreve que, na capitalização, ``se um deles fez o depósito e aos 85 anos o cidadão morreu, sabe qual o vencimento que ela fica? Zero!'', em referência à ausência de pensão. Também sustenta que, com contribuição reduzida, o trabalhador terminaria ``com exatamente um quarto, R\$750''.

  \item \textbf{Exemplos internacionais mostram que o sistema gerou problemas.}

  Ele diz que ``o sistema da capitalização já foi adotado no Chile, na Colômbia, no México e no Peru'' e que ``todos eles estão voltando atrás'', citando aposentadorias muito baixas e cobertura restrita. Como exemplo do Chile, menciona inclusive casos de ``suicídio entre os idosos'' e a ``média da aposentadoria no Chile'' como equivalente a ``meio salário mínimo''.

  \item \textbf{Ele usa dados e cenários para ilustrar o impacto sobre renda e o ganho do sistema financeiro.}

  No debate, cita uma estimativa segundo a qual ``a remuneração dos bancos consumiria R\$105.701,43'', o que equivaleria a mais de 62\% do patrimônio do empregado, deixando o trabalhador com menos. E conclui que ``um sistema sem empregador e com instituição financeira faturando um percentual será um fracasso, é condenar o trabalhador à miséria''.

  \item \textbf{Afirma que não há país que tenha aplicado esse modelo ``com recurso só do trabalhador'' e tido êxito.}

  Ele desafia: ``Se me apresentarem um único país que aplicou a tal de capitalização com o recurso só do trabalhador e que deu certo, eu mudo aqui o meu discurso''.
\end{enumerate}

\textbf{Exemplos usados por Paulo Paim}

\begin{itemize}
  \item Chile: aposentadorias baixas, menção a suicídio entre idosos e retorno parcial atrás do modelo.
  \item Chile, Colômbia, México e Peru: países citados como evidência de insucesso ou revisão do sistema.
  \item Cálculo ilustrativo em reais: estimativa de remuneração bancária consumindo grande parte do patrimônio.
  \item Cenário de pensão: caso em que o cônjuge remanescente ficaria com ``zero''.
  \item Desafio empírico: pedido de um único caso internacional bem-sucedido com recurso apenas do trabalhador.
\end{itemize}

\textbf{Limitações apontadas pela própria resposta}

\begin{itemize}
  \item O contexto não traz uma lista única e consolidada dos argumentos; eles aparecem distribuídos em vários trechos.
  \item A conclusão aparece dispersa em passagens distintas, e não em uma formulação final única.
\end{itemize}

\textbf{Referências retornadas pelo sistema}

\begin{itemize}
  \item Paulo Paim (PT/RS), 21/05/2019, \url{https://www25.senado.leg.br/web/atividade/pronunciamentos/-/p/texto/453686}
  \item Paulo Paim (PT/RS), 20/05/2019, \url{https://www25.senado.leg.br/web/atividade/pronunciamentos/-/p/texto/453766}
  \item Paulo Paim (PT/RS), 08/02/2019, \url{https://www25.senado.leg.br/web/atividade/pronunciamentos/-/p/texto/450736}
  \item Paulo Paim (PT/RS), 19/02/2019, \url{https://www25.senado.leg.br/web/atividade/pronunciamentos/-/p/texto/450992}
  \item Paulo Paim (PT/RS), 07/02/2019, \url{https://www25.senado.leg.br/web/atividade/pronunciamentos/-/p/texto/450672}
\end{itemize}

\subsection{Caso 2: reconhecimento de limite em desambiguação de autoria}

\textbf{Pergunta}

Quais propostas são realmente de Confúcio Moura e quais aparecem em discursos de outros senadores, mas poderiam ser confundidas com as dele?

\textbf{Resposta}

No contexto fornecido, \textbf{não é possível identificar com segurança} quais propostas são realmente de Confúcio Moura e quais aparecem em discursos de outros senadores, mas poderiam ser confundidas com as dele, porque o material recuperado traz apenas trechos e resumos de discursos, sem uma base completa que permita comparar cada proposta e distinguir, com segurança, autoria, citação ou mera concordância.

Dito isso, a resposta ainda consegue apontar, com base apenas no contexto recuperado, algumas linhas de atuação associadas a Confúcio Moura:

\begin{itemize}
  \item educação como prioridade, inclusive com a formulação de que ``a saída é a educação'';
  \item manifestação favorável à perenização do Fundeb e à ampliação de recursos;
  \item defesa de maior qualidade da educação básica e posição contrária à indicação política para diretorias escolares e secretarias municipais;
  \item intenção de proposta legislativa para melhorar condições de trabalho de catadores e recicladores.
\end{itemize}

Ao mesmo tempo, a resposta sinaliza corretamente o que \textbf{não foi possível confirmar}:

\begin{itemize}
  \item o contexto não traz uma separação completa entre ``propostas de Confúcio Moura'' e ``propostas de outros senadores'';
  \item alguns trechos aparecem em meio a debates ou falas de plenário, sem permitir distinguir com segurança autoria, apoio, comentário ou simples referência;
  \item por isso, o sistema evita afirmar além do que o material recuperado realmente sustenta.
\end{itemize}

\textbf{Por que este caso é relevante}

Este exemplo é metodologicamente importante porque mostra um tipo de bom comportamento que nem sempre aparece em demonstrações de sistemas RAG: em vez de forçar uma resposta completa sob incerteza, o sistema reconhece insuficiência de contexto, reduz o escopo da afirmação e preserva cautela na atribuição de autoria. Esse padrão é especialmente valioso em acervos parlamentares, nos quais temas semelhantes, vocabulário político recorrente e apartes em plenário podem induzir confusão entre autores.

\textbf{Referências retornadas pelo sistema}

\begin{itemize}
  \item Confúcio Moura (MDB/RO), 14/03/2019, \url{https://www25.senado.leg.br/web/atividade/pronunciamentos/-/p/texto/451323}
  \item Confúcio Moura (MDB/RO), 06/02/2020, \url{https://www25.senado.leg.br/web/atividade/pronunciamentos/-/p/texto/463649}
  \item Confúcio Moura (MDB/RO), 27/09/2019, \url{https://www25.senado.leg.br/web/atividade/pronunciamentos/-/p/texto/458160}
  \item Confúcio Moura (MDB/RO), 20/10/2020, \url{https://www25.senado.leg.br/web/atividade/pronunciamentos/-/p/texto/472033}
  \item Confúcio Moura (MDB/RO), 29/06/2021, \url{https://www25.senado.leg.br/web/atividade/pronunciamentos/-/p/texto/480774}
  \item Paulo Paim (PT/RS), 27/09/2019, \url{https://www25.senado.leg.br/web/atividade/pronunciamentos/-/p/texto/458158}
\end{itemize}

\subsection{Caso 3: pergunta multifocal e consolidação legislativa sob forte restrição de contexto}

\textbf{Pergunta}

Reconstrua, com base nos discursos recuperados, um plano legislativo consolidado para reduzir desigualdade no Brasil, indicando autores, medidas, cronograma, orçamento e consenso entre senadores.

\textbf{Resposta}

A resposta começa com um movimento metodologicamente correto: afirma que \textbf{não foi possível identificar, no conjunto de discursos recuperados, um plano legislativo único e consolidado} com autores, medidas, cronograma, orçamento e consenso explicitamente descritos dessa forma. Em vez de inventar uma síntese inexistente, o sistema reorganiza apenas o que o contexto realmente permite inferir, em cinco blocos: objetivo consolidado, medidas legislativas identificadas, cronograma, orçamento e consenso entre senadores.

\textbf{O que a resposta consegue fazer bem}

\begin{itemize}
  \item identifica objetivos recorrentes, como redução de desigualdades sociais, fortalecimento de políticas para população vulnerável, maior participação política de mulheres e busca de isonomia na distribuição orçamentária;
  \item recupera medidas legislativas concretas, citando, por exemplo, o PL n.\ 1.951/2021, a defesa da PEC n.\ 32/2022 e debates sobre orçamento secreto, emendas de bancada e limites de financiamento eleitoral;
  \item separa com clareza o que o contexto permite reconstruir e o que \textbf{não foi possível consolidar} com segurança, especialmente em relação a cronograma único, orçamento total e consenso formal entre senadores;
  \item evita apresentar como fato consumado um ``plano legislativo'' que, nos discursos recuperados, aparece apenas como conjunto fragmentário de propostas, críticas e agendas parciais.
\end{itemize}

\textbf{Por que este caso é mais desafiador}

Esta pergunta tensiona o sistema em múltiplas frentes ao mesmo tempo:

\begin{itemize}
  \item exige síntese entre vários autores e temas;
  \item pressupõe integração entre discursos que não necessariamente foram produzidos como parte de um mesmo plano;
  \item pede elementos estruturados --- cronograma, orçamento e consenso --- que raramente aparecem completos em falas parlamentares isoladas;
  \item aumenta o risco de o modelo ``costurar'' artificialmente uma política pública única a partir de trechos apenas parcialmente compatíveis.
\end{itemize}

\textbf{Limitações observadas}

\begin{itemize}
  \item embora a resposta seja cautelosa, ela já opera num nível de abstração maior do que os dois casos anteriores;
  \item a organização em blocos pode transmitir ao leitor a impressão de maior unidade programática do que aquela efetivamente sustentada pelos discursos;
  \item este é um bom exemplo de situação em que o sistema se comporta melhor ao reconhecer fragmentação do contexto do que ao tentar entregar uma resposta ``fechada''.
\end{itemize}

\textbf{Por que este caso é relevante}

Este terceiro exemplo ajuda a demonstrar um limite importante da solução: mesmo quando o sistema evita alucinação dura, perguntas amplas, multifocais e orientadas à consolidação normativa exigem supervisão humana mais forte. Em acervos parlamentares, a existência de temas comuns entre discursos não equivale, por si só, à existência de um plano legislativo único, consensual e plenamente estruturado.

\textbf{Referências retornadas pelo sistema}

\begin{itemize}
  \item Fabiano Contarato (REDE/ES), 10/04/2019, \url{https://www25.senado.leg.br/web/atividade/pronunciamentos/-/p/texto/452433}
  \item Orlando Silva (PCdoB/SP), 16/12/2022, \url{https://www25.senado.leg.br/web/atividade/pronunciamentos/-/p/texto/495399}
  \item Paulo Paim (PT/RS), 15/07/2021, \url{https://www25.senado.leg.br/web/atividade/pronunciamentos/-/p/texto/481744}
  \item Lasier Martins (PODEMOS/RS), 14/12/2021, \url{https://www25.senado.leg.br/web/atividade/pronunciamentos/-/p/texto/486225}
  \item Alessandro Vieira (PSDB/SE), 13/12/2022, \url{https://www25.senado.leg.br/web/atividade/pronunciamentos/-/p/texto/494690}
  \item Renan Calheiros (MDB/AL), 16/12/2022, \url{https://www25.senado.leg.br/web/atividade/pronunciamentos/-/p/texto/495400}
  \item Paulo Paim (PT/RS), 26/08/2019, \url{https://www25.senado.leg.br/web/atividade/pronunciamentos/-/p/texto/458186}
  \item Marcelo Castro (MDB/PI), 10/07/2019, \url{https://www25.senado.leg.br/web/atividade/pronunciamentos/-/p/texto/455943}
  \item Rodrigo Pacheco (PSD/MG), 20/12/2021, \url{https://www25.senado.leg.br/web/atividade/pronunciamentos/-/p/texto/486289}
\end{itemize}
